\section{Introduction}\label{sec:intro}

At the outset of a reverse-engineering project, the engineer must develop a sense of the program's high-level behavior in order to identify points of interest that warrant careful analysis.
When the target's input format is well understood, much of that triage reduces to identifying which instructions in the binary implement which features of the input language.
This paper presents a technique for doing exactly that.
The basic idea is to run the target on a large corpus of valid inputs while collecting control flow traces, identify corpus elements with similar parse trees, and associate the delta in their control flow with the delta between their parse trees.

Consider a proprietary binary that validates TLS certificates.
The analyst knows the input is an X.509 certificate and understands its grammar (version fields, subject names, key material, extensions) but has no source code.
\toolname takes the certificate grammar together with a set of structurally varied certificates, traces the binary under \QEMU, and produces a labeled control-flow graph in which each grammar element maps to the basic blocks that process it.
The analyst then begins reverse engineering with a semantically meaningful map of the parser rather than raw disassembly.

We present \toolname and evaluate it on four input formats spanning the Chomsky hierarchy~\cite{chomsky1956three}: a regular language (URI), a context-free language (JSON), a context-sensitive binary format with internal compression (HPACK), and a structured offset-driven binary format (ELF).
For each format we run \toolname on one or more popular open-source implementations, then validate by manually checking that each labeled basic block implements the grammar rule indicated by its label.

\textbf{Contributions.} Our paper makes the following contributions to the field of reverse engineering:
\begin{itemize}
    \item We introduce the first binary-only technique that transfers semantic labels from program inputs onto control-flow graph edges and nodes. This bridges the gap between input formats and program internals, assisting reverse engineering and parser understanding (Section~\ref{sec:design}).
    \item We implement our technique in \toolname, a Python tool that requires only \QEMU and the target binary. We release \toolname as open source and demonstrate it on four input formats from different language classes (Section~\ref{sec:implementation}).
    \textit{[BEN: confirm we can publish the public repo URL in the camera-ready. For dual-blind we leave it anonymous.]}
    \item We evaluate \toolname on real-world parsers (Curl, Apache APR, json-c, libelf, and nghttp2) and validate every label through manual inspection (Section~\ref{sec:evaluation}).
    \item Our work lays the groundwork for future applications in differential fuzzing, grammar mining, and vulnerability discovery (Section~\ref{sec:future}).
\end{itemize}
